% ============================================================
%  TD Chapitre 6 : Suites numeriques
% ============================================================
\clearpage
\chapter{Suites numériques}
\chaptermark{Suites numériques}

\begin{exercice}[monotonie de deux suites]
Étudier la monotonie des suites suivantes en calculant $u_{n+1} - u_{n}$ et en discutant son signe :
\begin{enumerate}
\item $u_{n} = 5 - 3n$ ;
\item $v_{n} = n^{2} + 2n$.
\end{enumerate}
\end{exercice}

\begin{exercice}[intérêts simples et valeur acquise]
Un capital de $250\,000$ F est placé à $6\,\%$ par an en intérêts simples, les intérêts produits chaque année ne sont pas replacés.
\begin{enumerate}
\item calculer le montant des intérêts produits chaque année, et expliquer pourquoi la valeur acquise suit une suite arithmétique dont on donnera le premier terme et la raison ;
\item calculer la valeur acquise après $5$ années ;
\item calculer la valeur acquise après $10$ années.
\end{enumerate}
\end{exercice}

\begin{exercice}[intérêts composés et comparaison]
Le même capital de $250\,000$ F est placé à $6\,\%$ par an en intérêts composés, chaque année les intérêts sont ajoutés au capital.
\begin{enumerate}
\item expliquer pourquoi la valeur acquise suit une suite géométrique dont on donnera le premier terme et la raison, puis écrire le terme général $V_{n}$ ;
\item calculer la valeur acquise après $5$ années, arrondie au franc près ;
\item calculer la valeur acquise après $10$ années, arrondie au franc près ;
\item comparer avec les résultats à intérêts simples des $5$ et $10$ ans, et dire quel mécanisme fait la différence.
\end{enumerate}
\end{exercice}

\begin{exercice}[amortissement linéaire]
Une machine de valeur $4\,000\,000$ F est amortie linéairement sur $8$ années : sa valeur nette comptable diminue chaque année d'une annuité constante.
\begin{enumerate}
\item calculer l'annuité d'amortissement annuelle, et écrire la valeur nette comptable $V_{n}$ après $n$ années comme le terme général d'une suite arithmétique dont on précisera la raison ;
\item calculer la valeur nette comptable après $5$ années ;
\item au bout de combien d'années la valeur nette comptable devient-elle inférieure ou égale à $1\,000\,000$ F ?
\end{enumerate}
\end{exercice}
